\documentclass[11pt]{article}
\usepackage[T1]{fontenc}
\usepackage{microtype}
\usepackage[letterpaper,margin=1.08in]{geometry}
\usepackage{amsmath,amssymb,amsthm,mathtools}
\usepackage[hidelinks]{hyperref}
\hypersetup{
  pdftitle={The Mathieu Property over Fields of Characteristic Zero: Classification and Ground-Field Obstructions},
  pdfauthor={Roy van Rijn}
}
\newtheorem{theorem}{Theorem}[section]
\newtheorem{lemma}[theorem]{Lemma}
\newtheorem{corollary}[theorem]{Corollary}
\newtheorem{proposition}[theorem]{Proposition}
\newcommand{\CC}{\mathbb C}
\newcommand{\RR}{\mathbb R}
\newcommand{\QQ}{\mathbb Q}
\newcommand{\Gm}{\mathbb G_m}
\newcommand{\SL}{\mathrm{SL}}
\newcommand{\SU}{\mathrm{SU}}
\newcommand{\SO}{\mathrm{SO}}
\newcommand{\PGL}{\mathrm{PGL}}
\newcommand{\CT}{\operatorname{CT}}
\newcommand{\tr}{\operatorname{tr}}
\newcommand{\Lie}{\operatorname{Lie}}
\newcommand{\dd}{\,\mathrm d}
\DeclareMathOperator{\Ad}{Ad}
\title{The Mathieu Property over Fields of Characteristic Zero:\\
\large Classification and Ground-Field Obstructions}
\author{Roy van Rijn}
\date{November 2026}
\begin{document}
\maketitle

\begin{abstract}
Let \(K\) be a field of characteristic zero and \(X\) an affine geometrically
homogeneous space of a group \(G\) with split reductive identity component.
The normalized Reynolds kernel on \(K[X]\) is Mathieu--Zhao if and only if
\([G^\circ,G^\circ]\) acts trivially. A single active split \(\SL_2\) already
gives a fixed marked counterexample for a connected reductive acting group.
Both results allow \(X(K)=\varnothing\). The proof transfers witnesses along
principal orbits and specializes marked-moment recurrences over the ground
field. We also classify connected reductive groups and compact adjoint
semisimple torsors over real closed fields, and representative functions
on compact Hausdorff homogeneous spaces. A long-root subgroup gives a
further obstruction when a split central torus acts alongside an active
compact factor of rank at least two.
\end{abstract}

\section{Introduction}

Let \(K\) be a field of characteristic zero. Let \(X\) be a smooth affine
\(K\)-variety with a geometrically transitive action of a linear algebraic
\(K\)-group \(G\) with reductive identity component. Thus \(X_{\overline K}\)
is a homogeneous quotient and its geometric stabilizers have reductive
identity component by the affine Matsushima criterion \cite{Matsushima}.
If \(x\in X(K)\), we may write \(X=G/H\) with \(H=G_x\).
The normalized Reynolds functional \(\Lambda_X:K[X]\to K\) projects onto
the invariant constants. A function \(P\) is \emph{pure} if
\(\Lambda_X(P^m)=0\) for every \(m\ge1\). It is \emph{safe} if for each
fixed \(Q\in K[X]\) one has \(\Lambda_X(QP^m)=0\) for all sufficiently
large \(m\). The kernel is Mathieu--Zhao precisely when every pure
function is safe. We call \((P,Q)\) a witness if \(P\) is pure and
\(\Lambda_X(QP^m)\ne0\) at infinitely many orders \(m\).

\begin{theorem}\label{thm:main}
Suppose \(G^\circ\) is split over \(K\). Then
\[
 \ker\Lambda_X\text{ is Mathieu--Zhao}
 \quad\Longleftrightarrow\quad
 [G^\circ,G^\circ]\text{ acts trivially on }X.
\]
For \(X=G/H\), the right-hand condition is
\([G^\circ,G^\circ]\subseteq H\). For connected \(G\), it also
characterizes equality of the pure set and the geometric nullcone on
\(K[X]\).
\end{theorem}

Mathieu's first conjecture posed the purity-to-safety implication for
representative functions on compact connected Lie groups \cite{Mathieu}.
His second asked whether pure functions on complex reductive spherical
pairs are unstable. Long classified the compact connected group case
and the complex spherical case \cite{Long,LongSpherical}.
The classification on arbitrary complex reductive homogeneous spaces
was proved in \cite{Complex}. Over the ground field the answer can change:
on the real quaternion norm-one group every nonzero real regular
function has positive squared trace, while its complexification is
\(\SL_2\) and admits a witness.

The negative direction of Theorem~\ref{thm:main} uses only one active
split \(\SL_2\) and works even when \(G\) is nonsplit. A general
stabilizer either supplies an adjoint covariant or determines a
normalizer quotient of \(\SL_2\). In the second case an invariant
denominator clears poles in the transferred witness. A recurrence
and Hilbert irreducibility select parameters so a single ground-field
multiplier detects infinitely many orders.

The absence of rational points also matters for compact groups.
Over a real closed field, a torsor for a compact adjoint semisimple
group has Mathieu--Zhao trace exactly when it has a rational point
(Theorem~\ref{thm:torsor}). For instance, the isomorphism torsor from
the quaternion division algebra to the split matrix algebra carries
commuting compact and split adjoint actions with the same normalized
trace.

We use the term \emph{geometric nullcone} for the points whose orbit
closure over an algebraic closure contains zero. For an individual
regular function this can be tested in any finite-dimensional rational
submodule containing it. For disconnected \(G\), geometric instability
is the same for \(G\) and \(G^\circ\), since the larger orbit closure is
a finite union of translates. An unstable function is safe. Over an
algebraic closure a contracting cocharacter gives strictly positive
weights on that function, eventually dominating the finitely many
weights of any fixed multiplier \cite{Kempf,LongSpherical}.
A nonzero marked subsequence therefore excludes geometric instability.

\section{Rank-one moment identities}\label{sec:rank}

Normalized Reynolds functionals commute with ground-field extension.
An equivariant morphism between transitive spaces preserves their
normalized functionals under pullback, by uniqueness. In particular,
on a quotient \(G/H\) the trace is the restriction of the normalized
left trace on \(K[G]\). We use these facts also for central isogenies
and products.

\begin{lemma}\label{lem:splitrank}
For split \(S=\SL_2\), the quotients by the split diagonal torus
normalizer and the two standard finite groups \(2O,2I\) admit
witnesses over \(K\).
\end{lemma}

\begin{proof}
Use the binary forms
\[
 f_4=X^2Y^2,\qquad
 f_8=X^8+14X^4Y^4+Y^8,\qquad
 f_{12}=XY(X^{10}+11X^5Y^5-Y^{10}).
\]
For a form of degree \(n=2b\) write \(a=b-1\). For the rows \(u,v\)
of a determinant-one matrix let \(U=f(u)\), \(F=f(v)\), and set
\[
 R=\frac{f(v+Fu)}{F},\qquad P_f=U^aR^b,\qquad Q_f=F^a.
\]
Taylor expansion in the added row shows \(R\) is polynomial. Every
subgroup fixing \(f\) fixes the pair under right translation.
For \(f_4,f_8,f_{12}\), the groups are respectively the torus
normalizer, \(2O\) and \(2I\). The pure moments vanish over \(\QQ\).
The quartic has nonzero marked moments at all positive orders.
For the two exceptional forms there are nonzero marked moments at
orders \(m\) for which \(28m-3\), respectively \(66m-5\), is prime
\cite[Theorem~1.2 and Lemma~5.1]{Complex}. In each progression the
residue is coprime to the modulus, so Dirichlet's theorem gives
infinitely many such orders.

\end{proof}

For use when a principal stabilizer is central, write a traceless
matrix as
\[
 Z=\begin{pmatrix}a/2&v\\u&-a/2\end{pmatrix},\quad r=a^2+4uv,
 \qquad P_{\rm ad}=v^2(1+au-u^3v)^4,\quad Q_{\rm ad}=u^2 .
\]
If \(R_S\) denotes the conjugation Reynolds projection, then
\begin{equation}\label{eq:adjoint}
 R_S(P_{\rm ad}^m)=0,\qquad
 R_S(Q_{\rm ad}P_{\rm ad}^m)=\beta_m r^{3m-1},\qquad
 \beta_m=\frac{2m}{3m+1}\frac{(2m)!(4m)!}{(6m)!}.
\end{equation}
For \(Z=ww^*-I/2\) with \(w=(x,y)^{\mathsf T}\) a unit vector,
\(u=\bar x y\), \(v=x\bar y\), \(a=|x|^2-|y|^2\), and
\[
                1+au-u^3v=(1+|x|^2u)(1-|y|^2u).
\]
Under the uniform unit-sphere measure, \(s=|x|^2\) is uniform
on \([0,1]\) and the angle of \(u\) is uniform. Put \(n=2m\).
Angular averaging selects degree \(k=n\) in \(u\) for the pure
moment and \(k=n-2\) for the marked one. Expanding the two factors
and integrating in \(s\) gives the common prefactor
\((2n)!^2/(2n+k+1)!\) times
\[
 \sum_{p+q=k}\frac{(-1)^q(n+p)!(n+q)!}
 {p!\,q!\,(2n-p)!\,(2n-q)!}.
\]
For \(k=n\), the factorial ratios cancel and the sum is zero.
For \(k=n-2\), it equals
\[
 \frac{1}{(n-2)!}\int_0^1\!\int_0^1
 x^ny^n(1-x)(1-y)(x-y)^{n-2}\dd y\,\dd x .
\]
Integrate on \(y=xt\), \(0<t<1\), and double by symmetry to
obtain the displayed factor \(\beta_m\).

At general radius, scaling \(Z\) by \(c\) multiplies \(r\) by \(c^2\).
A term of \(u\)-weight \(k\) in \((1+au-u^3v)^{4m}\) has total degree
\(2k\) in the matrix entries, with \(u,a,v\) of weights \(1,0,-1\).
The selected pure and marked terms thus have total degrees \(6m\)
and \(6m-2\). This proves \eqref{eq:adjoint} on the entire traceless
Hermitian real form. Zariski density gives the polynomial Reynolds
identities over every characteristic-zero field.

\section{Principal stabilizers}\label{sec:principal}

Suppose \(G\) is connected and let \(S\to G\) be a nontrivial rank-one
action with \(S_{\overline K}\simeq\SL_2\). We establish the principal
orbit statement for this geometric type, and then use splitness to
descend a particular normalizer quotient.

\begin{lemma}\label{lem:principal}
There is a dense \(S\)-invariant open \(U\subseteq X\) on which all
geometric \(S\)-orbits are closed and their stabilizers are conjugate.
The common stabilizer is proper and reductive.
\end{lemma}

\begin{proof}
Work first over \(\CC\) using an embedding of an algebraic closure of
a finitely generated field of definition. Choose a compact form \(C\)
of \(G_\CC\) containing the image of a compact form of \(S_\CC\).
After choosing a geometric base point, conjugate its stabilizer so
that a Zariski-dense compact subgroup of it lies in \(C\).
These choices are possible by compact-form conjugacy and containment
\cite[Theorems~3.7--3.8]{AdamsTaibi}. The conjugated quotient is
\(G_\CC\)-equivariantly isomorphic to the original. Compact conjugation
gives an antiholomorphic involution of the quotient and a Zariski-dense
compact \(C\)-orbit of fixed points.

Embed the quotient into a finite-dimensional \(S_\CC\)-module by
taking a stable space of algebra generators closed under the
involution. Average a Hermitian metric and choose it real on the
fixed real form. For any fixed real vector \(v\) and any
\(A\in\Lie(\SU_2)\), the vector \(Av\) is real. Hence
\(\operatorname{Re}\langle iAv,v\rangle=0\), which gives a closed
orbit by the Kempf--Ness criterion \cite{KempfNess}.
The compact orbit meets the intersection of the maximal orbit dimension
locus and the inverse image of the open principal Luna stratum
\cite[Part~III, \S3]{Luna}.
The closed orbit in each fiber over that stratum has the same
stabilizer type and thus maximal dimension. Any nonclosed orbit
there would have that closed orbit in its boundary and would have
strictly larger dimension. This proves the assertion on the inverse
image of the principal stratum. The geometric principal open
descends by uniqueness and base change of invariant algebras.
The closed orbits have reductive stabilizers by Matsushima
\cite{Matsushima}. If their stabilizer were all of \(S\), the action
would be trivial on a dense open.
\end{proof}

Suppose now \(S=\SL_{2,K}\) and let \(L\) represent the stabilizer
class over \(\overline K\). If \(L\) is noncentral, attach a group
\(M(L)\) as follows. When \(L^\circ\) is a torus, use its normalizer.
For a finite noncentral cyclic \(L\), use the normalizer of its
centralizer torus. For binary dihedral \(L\) of order \(4r\) with
\(r\ge3\), use the normalizer of the centralizer torus of the
characteristic rotation subgroup generated by elements of order
greater than four. For \(Q_8,2T,2O,2I\), use the full normalizer.
It is \(2O,2O,2O,2I\), respectively. Indeed \(Q_8\) and \(2T\)
are normal in \(2O\), the centralizer of a finite noncyclic subgroup
is the center, and the finite subgroup list rules out proper finite
overgroups of \(2O\) and \(2I\) \cite{Flicker}.

In all cases \(N_S(L)\subseteq M(L)\) and the assignment commutes
with conjugation and Galois. The group \(M(L)\) is conjugate to a
self-normalizing \(K\)-subgroup \(M_0\), namely the split diagonal
torus normalizer or the stabilizer of \(f_8\) or \(f_{12}\) in
Lemma~\ref{lem:splitrank}. The torus normalizer is self-normalizing
by its identity component. The two finite groups are self-normalizing
by maximality among finite groups. Conjugate \(L\) so \(M(L)=M_0\).
Over \(\overline K\) the rule
\[
             S_x=sLs^{-1}\ \longmapsto\ sM_0
\]
defines a rational equivariant map. Indeed, the image of
\((s,y)\mapsto(sy,sM_0)\) from \(S\times U^L\) is a constructible
graph with one geometric point over each point of \(U\).
For \(y\in U^L\), containment \(L\subseteq S_y\) is equality by
conjugacy, including equality of the numbers of components. The
unique dominating component of the graph closure is birational to
\(U\) by generic separability. For each Galois automorphism
\(\sigma\) write \(\sigma L=hLh^{-1}\). Since both attached groups
are \(M_0\), we have \(h\in M_0\). Replacing \(L\) by \(\sigma L\)
therefore leaves the graph unchanged. This descends the map
\(\phi:X\dashrightarrow S/M_0\) over \(K\) without a rational point of \(X\).

\subsection*{The two types of witness}

If \(L\) is central, restriction to a closed orbit shows that the
adjoint module occurs in \(K[X]\). Restriction is surjective over
the algebraic closure, the adjoint module occurs on \(S/L\), and
the intertwiner space commutes with base change. Hence there is
a \(K\)-defined nonzero equivariant map to the adjoint space.
There is also such a map over \(\CC\) taking Hermitian values on the
compact orbit. In fact the antilinear involution
\(
 Z\mapsto \tau_{\rm h}\circ Z\circ\tau_X
\)
on the space of equivariant maps uses the compact conjugation
\(\tau_X\) of the quotient and \(\tau_{\rm h}(M)=\bar M^{\mathsf T}\).
Its real fixed space spans the complex space. Choose a nonzero
fixed map. It has nonzero radius on an open subset of the compact
orbit. By \eqref{eq:adjoint}, therefore, all its positive marked
moments are nonzero. It can be written in a finite span of maps
defined over \(K\). This gives a finite-parameter family of pure
adjoint phases over \(K\), whose marked sequences are generically
nonzero.

If \(L\) is noncentral, pull back the witnesses \(P_0,Q_0\) of
Lemma~\ref{lem:splitrank} by \(\phi:X\dashrightarrow S/M_0\). Denote the two rational functions
by \(p,q\). Their translates span a finite-dimensional rational
submodule. Its common denominator ideal \(I\) in \(K[X]\) is nonzero
and \(S\)-stable. Over \(\CC\), choose a nonzero \(a\in I\). Write
\(\tau_X^*a\) for compact conjugation on functions.
Then \(d=R_S(a\,\tau_X^*a)\) belongs to \(I^S\).
On the compact orbit it is nonnegative, and it is positive somewhere
by density and the compact angular average. Hence
\(\Lambda_X(d^j)\ne0\) for every \(j\ge1\). The denominator ideal
and its invariant subspace commute with field extension. Indeed,
after choosing a finite basis \(a_j/b\) with a common denominator
\(b\), the ideal is the simultaneous kernel of multiplication
by the \(a_j\) into \(K[X]/(b)\). Therefore we may take a finite
parameter span of \(K\)-invariant denominators containing \(d\)
after field extension. For each member of that span the regular pair
\((dp,dq)\) is pure and has marked moments
\begin{equation}\label{eq:factor}
 \Lambda_X((dq)(dp)^m)
    =\Lambda_X(d^{m+1})\,\Lambda_{S/M_0}(Q_0P_0^m).
\end{equation}
Indeed, on the open of definition the \(S\)-projection gives constant
rank-one moments, times \(d^{m+1}\) in the marked case. The identities
extend by regularity, and \(\Lambda_X\circ R_S=\Lambda_X\).
The nonzero complex parameter above shows that the polynomial sequence
in \eqref{eq:factor} is not eventually zero.

The same denominator argument applies to any reductive subgroup of
\(G\), using a compatible compact form of \(G_\CC\). We will use this
observation for a rank-one group times a central torus.

\section{Arithmetic specialization}\label{sec:arithmetic}

\begin{lemma}\label{lem:specialize}
Let \(E/\QQ\) be finitely generated. Suppose
\(b_m(\boldsymbol u)\in E[\boldsymbol u]\) is not eventually zero over
\(E(\boldsymbol u)\) and satisfies a scalar polynomial recurrence
with nonzero first coefficient for all sufficiently large integer
indices. Then some specialization \(\boldsymbol u\in E^d\) keeps
infinitely many \(b_m\) nonzero.
\end{lemma}

\begin{proof}
If there are no parameters, there is nothing to specialize.
Write \(A(m,\boldsymbol u)\) for the first coefficient. Remove its
content in \(E[m]\) when tracking its zeros, and choose \(M\) beyond
the fixed integer zeros and the recurrence cutoff. Choose
\(m_0\ge M\) with \(b_{m_0}\ne0\) as a polynomial. On a line through
two suitably chosen \(E\)-points the reduced coefficient has no
constant root in \(m\), and \(b_{m_0}\) is still nonzero. To see this,
choose the first point with \(b_{m_0}\ne0\) and reduced coefficient
a nonzero polynomial in \(m\). Choose the second point to avoid
vanishing of the reduced coefficient at the finitely many roots
of the first evaluation. Each of these avoidance conditions is
possible since the content was removed. Denote the line parameter
by \(u\).

Factor the reduced coefficient in \(E(u)[m]\). Each linear factor
has a nonconstant root \(r_i(u)\). Choose a distinct rational
prime \(p_i\) and embed \(E\) into \(\QQ_{p_i}\) so that \(r_i(u)\)
has negative valuation on a nonempty \(p_i\)-adic open set. For
a finite pole \(\alpha\) of \(r_i\), embed \(E(\alpha)\) instead,
then take a small ball near the pole. A pole at infinity requires
only large \(p_i\)-adic absolute value. The embeddings exist for
infinitely many primes for any finitely generated field of
characteristic zero. Indeed, specialize a transcendence basis to
rational values where a minimal polynomial of a primitive element
has no coefficient poles or multiple roots. A root of the specialized
polynomial is algebraic over \(\QQ\), so it has a simple root in
\(\QQ_p\) for infinitely many \(p\), by Schur's prime-divisor theorem
(see \cite[Section~7]{Complex}) and Hensel's lemma. Perturb the rational parameters to algebraically
independent values sufficiently close in \(\QQ_p\), and lift the
simple root there to get the required embedding.

Choose a rational \(a\) in all the required local opens by weak
approximation. Put \(b=\prod_i p_i\), take rational \(c\ne0\)
sufficiently small at these primes, and choose a prime \(\ell\)
larger than the degree of \(A\) in \(m\). Make the substitution
\[
                         u=a+\frac{c}{1+b t^\ell}.
\]
For every \(t\in E\), the reciprocal denominator is integral at the
chosen places. For \(t\ne0\), the valuations of the two summands are
\(0\) and \(1+\ell v_{p_i}(t)\), distinct since \(v_{p_i}(t)\in\mathbb Z\).
Thus the local conditions always hold. The rational map has degree
\(\ell\), so \([E(t):E(u)]=\ell\). If \(\theta\) is a root of an
irreducible factor of degree \(e\ge2\), the degree of
\(E(t,\theta)/E(u)\) is divisible by \(\ell\) and \(e\), and is at
most \(\ell e\). Coprimality forces equality, preserving irreducibility.
If no local conditions are needed, take \(u=t\). The field \(E\) is
Hilbertian \cite[Chapter~13]{FriedJarden}. Specialize \(t\) to
preserve the degrees and irreducibility of the nonlinear factors.
Avoid poles of all coefficients and factors, zeros of the scalar
factor, and zeros of \(b_{m_0}\). Linear
roots have negative valuation at one assigned prime and hence
cannot be integers. The resulting first coefficient is nonzero
at every integer \(m\ge M\). Backward recurrence from an eventual
zero tail would give \(b_{m_0}=0\), a contradiction.
\end{proof}

\begin{lemma}\label{lem:shift}
Let \(X\) be affine and geometrically homogeneous under a connected
reductive group over a characteristic-zero field \(k\). For
\(F,B\in k[X]\) with \(F\ne0\), the sequence \(\Lambda_X(BF^m)\)
satisfies a scalar polynomial recurrence with nonzero first
coefficient at all sufficiently large integer indices.
\end{lemma}

\begin{proof}
A finite extension providing a point of \(X\) exists by taking the
residue field of a closed point. The orbit map pulls the functions
and the normalized trace back to the connected group. A recurrence
over the finite extension descends by expanding the coefficients
in a field basis and taking a component with nonzero first coefficient.
It therefore suffices to prove the assertion on the group.

The zero-dimensional case is constant. Otherwise take a nonzero
left-invariant volume form \(\omega\) on the group.
The frame \(D_i\) of infinitesimal left translations preserves it, and
\(\Lambda(D_i a)=0\). Every regular \((\dim G-1)\)-form is a sum of
\(a_i\iota_{D_i}\omega\), with differential \(\sum_i D_i(a_i)\omega\).
Thus the trace on top forms annihilates regular exact forms.

On the smooth affine open \(F\ne0\), over \(k(\nu)\), use the
connection \(\mathrm d+\nu\,\mathrm dF/F\wedge\). Its top de Rham
cohomology is finite-dimensional. Indeed, descend the data to a finitely
generated subfield, embed that field with the parameter \(\nu\) into
\(\CC\), and apply holonomic direct-image finiteness to this integrable
connection \cite[Theorem~3.2.3]{HTT}. On the affine open, cohomology of
the complex of global forms commutes with field extension.
This gives the algebraic Mellin shift for these moments
\cite[Section~1.3]{LoeserSabbah}. The classes
of \(B\omega,FB\omega,\ldots\) are linearly dependent. Drop any initial
zero coefficients and reindex the resulting recurrence. In a relation
among the forms, the primitive for the twisted exact term has bounded
\(F\)-denominators and finitely many poles in \(\nu\). At sufficiently
large integers \(\nu=m\), multiplication by \(F^m\) turns the
relation into an ordinary regular exact form. Applying the trace
and clearing parameter denominators proves the claim.
\end{proof}

\begin{theorem}\label{thm:active}
Over any characteristic-zero field, let \(X\) be affine and geometrically
homogeneous under a connected reductive group. If a homomorphism from
split \(\SL_2\) to the acting group induces a nontrivial action on \(X\),
then \(K[X]\) admits a witness, including when \(X(K)\) is empty.
\end{theorem}

\begin{proof}
Choose a finitely generated field of definition for the constructions
and their finite parameter spans. For central \(L\), use the polynomial
sequence \(\Lambda_X(r(Z)^{3m-1})\) and the nonzero factors
\(\beta_m\) in \eqref{eq:adjoint}. Its recurrence follows from
Lemma~\ref{lem:shift} with phase \(r^3\), multiplier \(r^2\), and an
index shift. For noncentral \(L\), use the entire product sequence
in \eqref{eq:factor}, which accounts for gaps in the rank-one marked
sequence. It is a mixed moment on \(X\times S/M_0\) with phase
\(d\otimes P_0\) and multiplier \(d\otimes Q_0\). In each case the
complex parameter from the preceding construction shows the sequence
is generically not eventually zero. Apply Lemma~\ref{lem:shift}
over the parameter field and clear denominators to obtain a polynomial
recurrence. Lemma~\ref{lem:specialize} gives one parameter value and
hence one fixed pair over the field of definition.
\end{proof}

\begin{proof}[Proof of Theorem~\ref{thm:main}]
The geometric components of \(X\) are the orbits of \(G^\circ\).
If \(D=[G^\circ,G^\circ]\) is active, choose a Galois orbit of
components with active \(D\)-action. It gives an open-and-closed
factor with finite separable constant field \(E/K\), a geometrically
integral homogeneous \(G^\circ_E\)-space \(Y\) over \(E\).
Some root subgroup of the split derived group acts nontrivially
on \(Y\), since their images generate the derived group.
Theorem~\ref{thm:active} supplies a pair on \(Y\).
On this factor the trace is \(r^{-1}\operatorname{Tr}_{E/K}\circ
\Lambda_Y\), where \(r\) counts all geometric components of \(X\).
This follows by base change and the equal-weight average of the
geometric component traces. Nondegeneracy of the field-trace
pairing shows that some member \(c\) of a fixed \(K\)-basis of \(E\)
satisfies \(\operatorname{Tr}_{E/K}(c\Lambda_Y(QP^m))\ne0\) at
infinitely many marked orders. Extend \(P,cQ\) by zero to get a
witness on \(X\).

If \(D\) acts trivially, the geometric components are homogeneous
spaces of a torus. The component average and finite-trace separation
of \cite[Section~7]{Complex}, followed by the Laurent theorem
\cite{DvK}, prove the Mathieu--Zhao property over the algebraic
closure, hence over \(K\). For connected \(G\), the Laurent
support half-space criterion also gives geometric instability
of every pure function. In the active case the marked witness
lies outside the nullcone by the weight bound.
\end{proof}

\section{Real closed fields and the nonsplit boundary}\label{sec:real}

\begin{theorem}\label{thm:realgroup}
Let \(K\) be real closed and \(G\) connected reductive. The normalized
trace on \(K[G]\) has Mathieu--Zhao kernel exactly when \(G\) is a torus
or \(G\) is \(K\)-anisotropic.
\end{theorem}

For the positive compact case, a nonzero \(f\) has nonzero squared
trace. One can embed a finitely generated field of definition of the
group and \(f\) into \(\RR\), preserving the negative definiteness of
the derived Killing form and the anisotropic central torus. Preserve the sign of the squared
trace as well if it is nonzero. For the first condition preserve the signs of the principal minors. The
anisotropic central torus over a real closed field is a product of
norm-one tori for \(K(i)/K\), so include this model isomorphism in
the field of definition. On the compact real identity group
\(\Lambda(f^2)>0\) by density of the real orbit and normalized Haar.
This also proves the positive assertion on every affine \(G/H\) with
such a connected compact acting group. The torus case follows from
\cite{DvK} after splitting.

We record the real specialization principle used here and below.
Given finitely many strict signs on elements of a finitely generated
ordered field, embed that field into \(\RR\) respecting those signs.
Write the field as a finite separable extension of
\(\QQ(t_1,\ldots,t_d)\). Express existence of the root, the signs,
nonzero denominators and the nonzero derivative by a finite formula.
It holds in a real closure, hence over \(\RR\) by quantifier
elimination \cite{Tarski}. The simple root persists on a real open
set of parameters. Choosing algebraically independent values in
that open gives the desired embedding. If \(d=0\), the simple real
root itself suffices. In particular a finite condition with
parameters in a real closed field follows if it holds at every real
embedding of a finitely generated field of definition. Otherwise
quantifier elimination would give finitely many signs certifying a
failure, which could be preserved by such an embedding.

For the first negative case, a noncompact derived group has a split
rank-one subgroup, so Theorem~\ref{thm:active} applies to \(G\).
The same argument applies to any \(G/H\) on which a noncompact almost
simple factor acts nontrivially. Indeed a nontrivial split
cocharacter in the derived group gives a rational unipotent and
hence a split rank-one homomorphism by rational Jacobson--Morozov
\cite{JM}. Over a real closed field the noncompact derived case can
equivalently be recognized by a positive derived Killing square.
Over \(\RR\), a vector in the noncompact summand of a Cartan
decomposition has self-adjoint adjoint endomorphism. A nonzero
eigenvector of nonzero eigenvalue is adjoint-nilpotent by weight
shift. The existence of a nonzero \(\mathfrak{sl}_2\)-triple is a
finite condition, so the real specialization principle gives
the same conclusion over \(K\).

The other negative case of Theorem~\ref{thm:realgroup} has nontrivial
compact derived group and split central direction. Here is a
ground-field construction that also works over arbitrary fields.

\begin{proposition}\label{prop:central}
Suppose \(G\) is connected reductive, it has a nontrivial \(K\)-character
\(\chi\), and its derived group contains a \(K\)-subgroup \(S\) of
geometric type \(A_1\). Then there is a fixed marked counterexample
on \(K[G]\).
\end{proposition}

\begin{proof}
Write \(q\) for the negative of the derived Killing form on
\(\mathfrak s=\Lie(S)\), with polarization \(\langle\,,\,\rangle_q\).
Its restriction is nondegenerate, as one sees after splitting \(S\)
by evaluating the sum of squared weights on a coroot. Let \(p\) be
the Killing-orthogonal projection to \(\mathfrak s\), choose
\(\xi\in\mathfrak s(K)\) with \(q(\xi)\ne0\), and set
\[
 z=p(\Ad(g)\xi),\quad A=q(z),\quad
 B=\frac{\langle\xi,z\rangle_q^2}{q(\xi)},\quad t=\chi(g).
\]
Use
\begin{equation}\label{eq:angular}
 P=t^{-1}(1+t)\{A-(1+t)^2B\}^{2},\qquad Q=t .
\end{equation}
To check moments, use a model embedded in \(\CC\). By conjugating
the compact form of \(S\) we may write \(\xi=c\xi_0\) with \(\xi_0\)
compact and \(c\ne0\). This uses transitivity of the adjoint action
on each nonzero quadric level. Extend that compact form to the
derived group. Replacing \(\xi_0\) by \(c\xi_0\) scales \(P\) by
\(c^4\). For \(\xi_0\), the law of \(z\) under compact derived Haar
measure is rotation-invariant in the three-dimensional compact
adjoint space. On each nonzero radius \(B/A=w^2\) with \(w\)
uniform on \([0,1]\). The central isogeny from the product of the
connected center and the derived group separates the trace.
The variable \(t\) is a nontrivial central character, so central
averaging extracts its constant term.

Put \(F_m(s)=\int_0^s(1-w^2)^{2m}\dd w\). The angular average of
\(P^m/A^{2m}\) is
\[
                  t^{-m}(1+t)^{m-1}F_m(1+t).
\]
Since \(F_m(1+t)-F_m(1)\) is divisible by \(t^{2m+1}\), its constant
term is zero. Its constant term after multiplication by \(t\)
is \(F_m(1)>0\). The remaining radial integral of \(A^{2m}\)
is positive, since \(z\) is nonzero at the identity. This proves
the two moment statements after rescaling and hence over \(K\).
The same pair descends to \(G/H\) if \(H\) fixes \(\xi\) and
\(\chi|_H=1\).
\end{proof}

For a nontrivial compact derived group over a real closed field,
a maximal torus in it is anisotropic. Conjugation acts by \(-1\)
on its character lattice, so an opposite pair of roots generates
a rank-one subgroup over \(K\). If the center has a split direction,
there is a nontrivial \(K\)-character, and Proposition~\ref{prop:central}
applies. This completes Theorem~\ref{thm:realgroup}.

\subsection{Active long-root subgroups}

\begin{theorem}\label{thm:longroot}
Let \(X\) be affine and geometrically homogeneous under a connected
reductive group \(G\) over a characteristic-zero field \(K\).
Suppose a split central torus \(T\) acts nontrivially. If \(G\) contains
a \(K\)-subgroup \(S\) which geometrically is a long-root \(\SL_2\)
in an active almost simple factor of rank at least two, then \(K[X]\)
admits a witness. Over a real closed field such a subgroup exists
in every active compact almost simple factor of rank at least two.
\end{theorem}

\begin{proof}
Choose a geometric simple system in which the long root \(\beta\)
is simple.
For a neighboring simple root \(\alpha\), the Cartan integer
\(\langle\alpha,\beta^\vee\rangle=-1\), since \(\beta\) is long.
The element \(\beta^\vee(-1)\) is therefore noncentral in the
almost simple factor. It acts nontrivially since the factor acts
nontrivially and its kernel is central. If this central element
of \(S\) belonged to a principal stabilizer, it would fix the
entire principal open of Lemma~\ref{lem:principal}. Thus a finite
principal stabilizer omits the unique involution in \(S\). It has
odd order and is cyclic. A proper positive-dimensional reductive
principal stabilizer has a torus as identity component.

We first choose a nonzero \(S\)-invariant \(t\in K[X]\) of
nontrivial \(T\)-weight. In fact every irreducible module of the
almost simple factor over \(\overline K\) has a vector fixed by the long-root
subgroup. To see this, pass to its simply connected cover. If
there are at least two long simple roots, for each fundamental
weight choose a long simple root whose coroot pairs to zero with
that weight. The associated rank-one subgroup fixes the highest
vector. All long roots are conjugate. The root systems with only
one long simple root are \(C_r\) and \(G_2\), with \(B_2=C_2\).
For \(G_2\) the fundamental module for the long simple root is
the adjoint module, with a fixed vector in the kernel of that
root on the Cartan subalgebra. For \(C_r\) it suffices to handle
the last fundamental module, the primitive exterior power of
degree \(r\) of the symplectic module. With the long-root subgroup
on \(e_1,f_1\), a lowering operator for \(-\epsilon_1-\epsilon_r\)
applied to \(e_1\wedge\cdots\wedge e_r\) gives, up to a scalar,
\[
 (e_1\wedge f_1-e_r\wedge f_r)\wedge e_2\wedge\cdots\wedge e_{r-1}.
\]
This vector is fixed by the long-root subgroup. For sums of
fundamental weights, multiply nonzero left \(S\)-invariants in
the spaces of right highest-weight matrix coefficients in the
coordinate ring of the simply connected group. Their product
is nonzero and lies in the space for the sum of the weights.
The matrix coefficient decomposition identifies that space
with the dual irreducible module, and duality preserves
the existence of invariants. This proves the assertion for
every irreducible. Each nonzero \(T\)-weight space in \(K[X]\)
therefore has \(S\)-invariants by base change, and the active
split torus supplies one of nontrivial weight.

Write \(q\) for the negative Killing form on \(\mathfrak s=\Lie(S)\),
and choose \(\xi\in\mathfrak s(K)\) with \(q(\xi)\ne0\).
If the principal stabilizer is nontrivial, it determines a
nonisotropic line in \(\mathfrak s\). This is the Lie algebra of
its torus centralizer in the cyclic case or of its identity
component in the positive-dimensional case. The line is rational
in the point by the graph argument of Section~\ref{sec:principal}
over \(\overline K\). Since it depends only on the stabilizer,
it descends and is \(T\)-invariant. Put \(A=1\) and
\[
                  B=\frac{\langle z,\xi\rangle_q^2}{q(z)q(\xi)}
\]
for a representative \(z\) of the line. If the principal stabilizer
is trivial, take a nonzero regular adjoint covariant as in
Section~\ref{sec:principal}, with coordinate functions of one
\(T\)-weight. Divide by a nonzero \(S\)-invariant function of that
weight and clear its denominator by an \(S\times T\)-invariant
element using the compact square argument. This gives a regular
\(T\)-invariant adjoint covariant \(Z\ne0\). Set \(A=q(Z)\) and
\(B=\langle Z,\xi\rangle_q^2/q(\xi)\).

In both cases \(R_S(B^j)=A^j/(2j+1)\) on the open of definition.
For a compact real axis this is the spherical average of the
\(2j\)-th power of the direction cosine. For vectors the same
averaging formula is polynomial and hence holds throughout the
complexified adjoint space. For lines it follows on their
homogeneous space from the compact orbit. Define
\[
 P=d\,t^{-1}(1+t)\bigl(A-(1+t)^2B\bigr)^2,\qquad Q=t,
\]
where \(d\) is an \(S\times T\)-invariant denominator. The
expressions to be cleared are finite under \(S\times T\):
the axis terms are matrix coefficients of the rank-one group,
and \(t^{-1}\) is a weight vector. In the covariant case a
denominator for \(t^{-1}\) alone suffices as \(Z\) varies.
Taking the constant term in \(t\), after the \(S\)-average,
uses the same integral as in \eqref{eq:angular}. For all \(m\ge1\),
\[
        \Lambda_X(P^m)=0,\qquad
        \Lambda_X(tP^m)=F_m(1)\Lambda_X((dA^2)^m).
\]
In particular these identities are valid for finite parameter
spans of invariant denominators and, when needed, covariants.

Over \(\CC\), the denominator can be chosen nonnegative and positive
somewhere on the compact orbit by the compact square construction.
In the covariant case the regular \(T\)-invariant adjoint maps
are stable under compact conjugations, so choose a nonzero map
with compact real values on that orbit. The squared radius is
then positive on a dense open subset of the compact orbit.
Consequently \(\Lambda_X((dA^2)^m)\) is positive for every
positive order at those complex parameters. The coaction equations
and the common denominator ideal commute with base change.
Choose finite parameter spans over a finitely generated field of
definition whose complexifications contain these choices.
Lemmas~\ref{lem:shift} and \ref{lem:specialize} give one ground-field
parameter with infinitely many nonzero marked moments.

Over a real closed field, a maximal torus of a compact factor
is anisotropic and conjugation acts on its roots by inversion.
The subgroup for the opposite long-root pair descends to \(K\).
This proves the last assertion.
\end{proof}

\subsection{Compact adjoint torsors}

\begin{theorem}\label{thm:torsor}
Let \(K\) be real closed and \(C\) a connected adjoint semisimple
group of compact type. For a \(C\)-torsor \(X\), the normalized
trace on \(K[X]\) has Mathieu--Zhao kernel if and only if
\(X(K)\ne\varnothing\).
\end{theorem}

\begin{proof}
If \(X\) has a point, identify it with \(C\) and use strict square
positivity. In general write \(X\) as a right torsor. After choosing
a point over \(K(i)\), the descent involution has the form
\(g\mapsto c\,\sigma(g)\), with \(c\,\sigma(c)=1\) and \(\sigma\)
the original conjugation of \(C\). The commuting group of left
translations descends to the inner form \(C_c\) with conjugation
\(\operatorname{Int}(c)\sigma\). Its action is geometrically
transitive. The two normalized traces coincide: each commuting
action preserves the other's normalized trace by uniqueness,
and each action has only constant invariants.

If \(C_c\) is compact, its conjugation is conjugate to \(\sigma\)
by an inner automorphism \(\operatorname{Int}(a)\). Over \(\RR\)
this is conjugacy of compact forms \cite[Theorem~3.8]{AdamsTaibi}.
Over \(K\) it follows by the real specialization principle above.
Negative definiteness of the Killing forms and the existence
of \(a\) satisfying the conjugacy identity are finite algebraic
conditions in a presentation of the group and its involutions.
Adjointness then gives \(c=a\sigma(a)^{-1}\), so the point with
coordinate \(a\) descends to \(K\).
Thus if \(X(K)\) is empty, \(C_c\) has positive split rank.
Its action gives a witness by Theorem~\ref{thm:active}, and the
same witness applies to the trace for \(C\).
\end{proof}

The point criterion also holds for compact semisimple groups with
odd-order center. A cocycle that is trivial in the adjoint
cohomology set comes from the center. Conjugation fixes every
geometric point of the center, since the center lies in a compact
maximal torus. A central cocycle therefore has order dividing two
and is trivial when the center has odd order. A nontrivial adjoint
class gives a noncompact commuting inner form by the proof above.

For example, let \(D\) be a quaternion division algebra over any
characteristic-zero field and \(X=\operatorname{Isom}_{K\text{-alg}}
(D,M_2(K))\). A nonzero reduced-trace-zero \(\xi\in D\) gives the
split adjoint covariant \(Z(\phi)=\phi(\xi)\). Its radius is
\(r(Z)=4\xi^2\in K^\times\). Thus \eqref{eq:adjoint} gives an
explicit pair with nonzero marked moments at every positive order.
By the commuting-action argument this is also a witness for the
action of \(\PGL_1(D)\).

\subsection{Open ground-field cases}

The split obstruction does not classify individual anisotropic
nonsplit spaces over general fields. A basic arithmetic example
is the norm-one group of the quaternion division algebra
\((2,3)_{\QQ}\), which has positive real split rank but no
\(\QQ\)-split rank-one subgroup.

Even over a real closed field, the connected group classification
and the adjoint torsor theorem leave general homogeneous spaces
open. Two test cases are
\[
                (\SU_2/2O)\times\Gm,\qquad
                (\SU_2/2I)\times\Gm ,
\]
where \(\SU_2\) denotes the compact algebraic form. A nontrivial
finite cyclic stabilizer determines an axis and leads to the
angular identity. The binary octahedral and icosahedral stabilizers
have no invariant axis. The complex witnesses for their split
quotients do not give witnesses over the real ground field.

\section{Compact homogeneous spaces}

\begin{theorem}\label{thm:compact}
Let \(C\) be a compact Hausdorff group and \(L\) a closed subgroup.
On the complex representative functions of \(C/L\), normalized
invariant integration has Mathieu--Zhao kernel exactly when
\(\overline{[C^\circ,C^\circ]}\subseteq L\). In particular on \(C\)
it holds exactly when \(C^\circ\) is abelian.
\end{theorem}

\begin{proof}
In the compact Lie case complexify the pair. A faithful unitary
model has reductive Zariski closure, and its regular functions
restrict to all representative functions. The matrix-entry algebra
is conjugation-stable on a unitary group, so it is dense by
Stone--Weierstrass. Projecting onto each finite-dimensional
Peter--Weyl isotype proves the asserted equality. The Zariski
closure of \(L\) meets the compact Lie group in \(L\): average
matrix-entry polynomials over \(L\) to separate distinct right
cosets. Right invariance and normalized integration therefore agree
with the complex homogeneous-space constructions. The identity
component of a compact Lie group is Zariski dense in the corresponding
algebraic identity component, and the same holds for its derived
subgroup. Apply the complex classification.

For a general compact Hausdorff group, any two representative functions
factor through a single compact Lie quotient \(J\), using the direct
sum of their finite-dimensional representations. If the functions
are right \(L\)-invariant, they descend to \(J/L_J\) with
\(L_J\) the image of \(L\). The image of \(C^\circ\) is \(J^\circ\)
because the quotient by that image is a quotient of the profinite
group \(C/C^\circ\). Haar measure pushes forward, so the containment
criterion suffices by the compact Lie case. Conversely, an element
of \([C^\circ,C^\circ]\) outside \(L\) can be separated from \(L\)
by a finite sum of representations, using point separation and
compactness of \(L\). The resulting compact Lie quotient provides
a counterexample which pulls back \cite{Fausk}.
\end{proof}

If one defines instability using the Zariski closure of the compact
orbit in the finite-dimensional complex orbit span, the corresponding
nullcone criterion holds precisely when the lattice criterion of
\cite[Section~7]{Complex} holds on every compact Lie quotient pair.
Both the moments and the orbit of an individual representative
function factor through such a quotient. In particular for connected
\(C\), the same containment criterion governs both properties.

\section{Universal Reynolds projections and toral nullcones}

\begin{theorem}\label{thm:universal}
For a linear algebraic group \(G\) with reductive identity component
over \(K\), the following are equivalent.
\textup{(i)} \(G^\circ\) is a torus.
\textup{(ii)} The Reynolds projection has Mathieu--Zhao kernel on every
commutative rational \(G\)-algebra.
\textup{(iii)} For every such algebra and every \(f\) with zero Reynolds
projection of each positive power, each fixed geometric isotypic
component of \(f^m\) vanishes eventually.
\end{theorem}

\begin{proof}
Assume \(G^\circ=T\) is a torus. Work over the algebraic closure
in the finitely generated stable subalgebra \(B\) generated by a
finite-dimensional submodule containing \(f\). Choose component
representatives \(g_1=1,\ldots,g_s\), and write \(L_i\) for the
Laurent polynomial recording the \(T\)-weights of \(g_i f\) with
coefficients in \(B\). Since the component quotient acts by averaging
the \(T\)-projection, \(\sum_i\CT L_i^m=0\) for every \(m\ge1\).

Take a primary decomposition \(0=\bigcap_j\mathfrak q_j\) in \(B\).
Put \(\mathfrak p_j=\sqrt{\mathfrak q_j}\), and choose \(e_j\) with
\(\mathfrak p_j^{e_j}\subseteq\mathfrak q_j\).
In the fraction field of \(B/\mathfrak p_j\), finite-trace separation
\cite[Lemma~7.1]{Complex} gives pure constant terms for each Laurent
polynomial separately. Write \(L_1=U+V\) by putting the coefficients
in \(\mathfrak p_j\) in \(V\). If \(U\ne0\), the Laurent theorem gives
a functional strictly positive on its support. Terms of \((U+V)^m\)
with fewer than \(e_j\) factors of \(V\) then eventually cannot
reach a prescribed weight. All other terms have coefficients in
\(\mathfrak q_j\). The case \(U=0\) is immediate. Thus each prescribed
weight of \(f^m\) is eventually zero by the intersection of the
primary components. This proves both conclusions: each fixed
multiplier and each fixed geometric simple module has only finitely
many \(T\)-weights.

If \(G^\circ\) is non-toral, Theorem~\ref{thm:main} after algebraic
closure gives a pair on the group that is defined over a finite
extension \(E/K\). Regard \(E[G_E]\) as a rational \(G\)-algebra
over \(K\). Its invariants are \(E\), and its Reynolds projection
is the normalized trace with values in \(E\). The pair disproves
the universal kernel assertion. It also disproves eventual isotypic
vanishing, since the trace of a product pairs only dual types and
a fixed multiplier has finitely many types.
\end{proof}

Finally suppose \(D=[G^\circ,G^\circ]\) acts trivially on \(X\),
with \(G^\circ\) possibly nonsplit. Let \(T=G^\circ/D\), \(M=X^*(T_{\overline K})\), and let
\(\Gamma\) be a finite Galois group acting on \(M\) and the geometric
components. For a representative \(i\) of each component orbit
\(\alpha\), let \(\Gamma_i\) be its stabilizer and \(M_i\subset M\)
the weight lattice of that component. Put
\(V_i=(M_i\otimes\QQ)^{\Gamma_i}\), \(V=(M\otimes\QQ)^\Gamma\).
The ground-field pure set equals the geometric nullcone on \(K[X]\)
exactly when
\begin{equation}\label{eq:transfer}
 \bigoplus_\alpha V_i\longrightarrow V,\qquad
 (v_i)\longmapsto\sum_\alpha|\Gamma|^{-1}
                                  \sum_{\gamma\in\Gamma}\gamma v_i
\end{equation}
is injective. In fact purity separates componentwise over
\(\overline K\). The convex hull of each component support avoids
zero, and this remains true for its stabilizer-averaged weights
because the support is \(\Gamma_i\)-stable. If \eqref{eq:transfer}
is injective, strictly positive functionals on these averaged
supports combine into one invariant functional on \(M\otimes\QQ\).
It is positive on the original weights and yields a contracting
cocharacter after clearing denominators and lifting through the
connected center. Conversely, a nonzero integral relation in
\eqref{eq:transfer} gives one nonzero weight on each of the
participating Galois orbits of components. The corresponding
direct sum of character lines descends, and a nonzero \(K\)-vector
on each participating orbit is nonzero on every geometric
component of that orbit. This produces a pure function whose
global support has zero in its convex hull. It is therefore
outside the geometric nullcone. For split \(T\) this uses one
embedded character lattice per Galois orbit.

\subsection*{Two ground-field transfers}

Let \(E/K\) be finite separable. A fixed pair on \(Y=G/H\) over \(E\)
gives a fixed pair on its Weil restriction to \(K\). After scalar
extension to \(\overline K\), that restriction is the product of
the conjugates \(Y^\sigma\). The products of the conjugate phases
and multipliers are defined over \(K\), and their marked moment
is \(N_{E/K}(\Lambda_Y(QP^m))\). Purity and the precise nonzero
marked orders therefore descend to this restricted space.

Disconnected forms can introduce an additional field issue. For
a real closed \(K\), let \(C\) be the compact quaternion norm-one
group and take \(G=C\times\mu_3\). The two nonreal components are
represented by the \(K(i)[C_{K(i)}]\) factor of \(K[G]\).
Use on it a complex rank-one pair with rational marked moments
and extend by zero on the other factor. The component average
gives a bad pair over \(K\) with marked moments scaled by \(2/3\).
Thus the connected group classification over real closed fields
does not extend by inspecting the real identity component alone.

\paragraph{AI assistance and review.}
ChatGPT (OpenAI) assisted with mathematical arguments, literature checks,
exposition, proof review and verification code.

\begin{thebibliography}{99}
\small
\setlength{\itemsep}{2pt}
\setlength{\parsep}{0pt}
\bibitem{Mathieu}
O. Mathieu, Some conjectures about invariant theory and their applications,
\emph{S\'emin. Congr.} 2 (1997), 263--279.
\href{https://smf.emath.fr/sites/default/files/2018-11/smf_sem-cong_2_263-279__sample.pdf}{Publisher excerpt}.

\bibitem{Long}
C. D. Long, The Mathieu property for compact connected Lie groups,
\href{https://arxiv.org/abs/2609.16178}{arXiv:2609.16178}.

\bibitem{LongSpherical}
C. D. Long, A classification of Mathieu's second conjecture for
reductive spherical pairs, September 2026,
\href{https://github.com/long-mathematics/mathieu-second-conjecture-spherical-pairs/blob/e9edac262b4b64eb68feb0e6e80496cefe4a1dee/mathieu_second_conjecture_spherical_pairs.tex}{versioned source e9edac2}.

\bibitem{Complex}
R. van Rijn, Counterexamples to the Mathieu property: from explicit
quotient constructions to a general classification, 2026.
\href{https://hexagonmath.org/2611.00556v1}{Hexagon 2611.00556v1},
\href{https://doi.org/10.5281/zenodo.24310836}{doi:10.5281/zenodo.24310836}.

\bibitem{DvK}
J. J. Duistermaat and W. van der Kallen, Constant terms in powers of
a Laurent polynomial, \emph{Indag. Math.} 9 (1998), 221--231.
\href{https://doi.org/10.1016/S0019-3577(98)80020-7}{DOI}.

\bibitem{Flicker}
Y. Z. Flicker, Finite subgroups of \(\SL(2,F)\) and automorphy,
\emph{Manuscripta Math.} 167 (2022), 213--243.
\href{https://doi.org/10.1007/s00229-021-01274-x}{DOI}.

\bibitem{Luna}
D. Luna, Slices \'etales, \emph{M\'em. Soc. Math. France} 33 (1973),
81--105. \href{https://numdam.org/articles/10.24033/msmf.110/}{Numdam}.

\bibitem{KempfNess}
G. Kempf and L. Ness, The length of vectors in representation spaces,
\emph{Lecture Notes in Math.} 732 (1979), 233--243.
\href{https://doi.org/10.1007/BFb0066647}{DOI}.

\bibitem{Kempf}
G. Kempf, Instability in invariant theory, \emph{Ann. of Math.}
108 (1978), 299--316. \href{https://arxiv.org/abs/1807.02890}{Transcription}.

\bibitem{AdamsTaibi}
J. Adams and O. Ta\"ibi, Galois and Cartan cohomology of real groups,
\emph{Duke Math. J.} 167 (2018), 1059--1097.
\href{https://arxiv.org/abs/1611.05956}{arXiv:1611.05956}.

\bibitem{Matsushima}
I. V. Arzhantsev, Invariant ideals and Matsushima's criterion,
\href{https://arxiv.org/abs/math/0506430}{arXiv:math/0506430}.

\bibitem{FriedJarden}
M. D. Fried and M. Jarden, \emph{Field Arithmetic}, fourth ed.,
Springer, 2023. \href{https://doi.org/10.1007/978-3-031-28020-7}{DOI}.

\bibitem{LoeserSabbah}
F. Loeser and C. Sabbah, \'Equations aux diff\'erences finies et
d\'eterminants d'int\'egrales de fonctions multiformes,
\emph{Comment. Math. Helv.} 66 (1991), 458--503.
\href{https://perso.pages.math.cnrs.fr/users/claude.sabbah/articles/loeser-sabbah91.pdf}{PDF}.

\bibitem{Tarski}
C. Cohen and A. Mahboubi, Formal proofs in real algebraic geometry:
from ordered fields to quantifier elimination,
\emph{Log. Methods Comput. Sci.} 8 (2012).
\href{https://arxiv.org/abs/1201.3731}{arXiv:1201.3731}.

\bibitem{JM}
R. Appenzeller, Semisimple algebraic groups over real closed fields,
\href{https://arxiv.org/abs/2601.07732}{arXiv:2601.07732}, \S5.8.

\bibitem{HTT}
R. Hotta, K. Takeuchi and T. Tanisaki,
\emph{D-Modules, Perverse Sheaves, and Representation Theory},
Progress in Mathematics 236, Birkh\"auser, 2008.
\href{https://doi.org/10.1007/978-0-8176-4523-6}{DOI}.

\bibitem{Fausk}
H. Fausk, Equivariant homotopy theory for pro-spectra,
\emph{Geom. Topol.} 12 (2008), 103--176, Appendix A.
\end{thebibliography}
\end{document}
